\documentclass[11pt, a4paper]{article}

\usepackage[T1]{fontenc}
\usepackage[utf8]{inputenc}
\usepackage{tgpagella}
\usepackage[spanish, es-noshorthands]{babel}
\usepackage{graphicx}
\usepackage[most]{tcolorbox}
\usepackage[american]{circuitikz}
\usepackage{tikz}
\usetikzlibrary{shapes.geometric, arrows.meta, positioning}
\usepackage[margin=2.5cm]{geometry}
\usepackage{fancyhdr}

\pagestyle{fancy}
\fancyhf{}
\setlength{\headheight}{16pt}
\lhead{\textbf{UCB Sede Tarija} | Guía Docente}
\chead{Defensa Individual Diferida (Hito 1)}
\rhead{Circuitos Electrónicos II}
\renewcommand{\headrulewidth}{0.4pt}

\definecolor{ucbblue}{RGB}{0, 51, 102}

\begin{document}

\begin{tcolorbox}[colback=ucbblue!5!white, colframe=ucbblue, title=\textbf{Guía de Entrevista: Defensa Individual Hito 1}]
    Esta guía convierte la defensa grupal en una comprobación individual de 10-15 min, manteniendo el mismo rigor técnico[cite: 5]. \textbf{No debe proporcionarse a los estudiantes.[cite: 5]}
\end{tcolorbox}

\section*{Arquitectura del Sistema (Mapa Visual)[cite: 5]}
\begin{center}
    \begin{tikzpicture}[
        box/.style={rectangle, draw=ucbblue, thick, fill=white, text width=2.5cm, align=center, rounded corners, minimum height=1cm, font=\footnotesize},
        arrow/.style={-Stealth, thick, ucbblue}
    ]
        \node[box] (ntc) {1. Sensor NTC};
        \node[box, right=0.8cm of ntc] (wheat) {2. Puente Wheatstone AC};
        \node[box, right=0.8cm of wheat] (twin) {3. Twin-T Notch};
        \node[box, right=0.8cm of twin] (out) {4. Salida del Sistema};
        
        \node[below=1cm of wheat, font=\footnotesize\itshape, color=gray] (lt) {LTSpice / Osciloscopio};
        \node[below=1cm of out, font=\footnotesize\itshape, color=gray] (py) {CSV / Python};

        \draw[arrow] (ntc) -- (wheat);
        \draw[arrow] (wheat) -- node[above]{\footnotesize $V_{dif}$} (twin);
        \draw[arrow] (twin) -- (out);
        
        \draw[dashed, gray] (wheat) -- (lt);
        \draw[dashed, gray] (out) -- (py);
    \end{tikzpicture}
\end{center}

\begin{center}
    \begin{circuitikz}[scale=0.9, transform shape, thick]
        \node[above] at (1.5,2.5) {\textbf{Modelo Conceptual del Efecto de Carga (Loading)[cite: 5]}};
        \draw (0,0) to[V, v=$V_{th}$] (0,2) to[R, l=$Z_{out}$] (2,2) coordinate(out) to[short, -o] (3,2);
        \draw (0,0) -- (2,0) to[short, -o] (3,0);
        \draw (3,2) to[R, l=$Z_{in}$] (3,0);
        \node[above] at (3,2) {$V_L = V_{th} \frac{Z_{in}}{Z_{out}+Z_{in}}$[cite: 5]};
    \end{circuitikz}
\end{center}

\section*{Secuencia de la Entrevista (15 mins max)[cite: 5]}
\begin{enumerate}
    \item \textbf{0–2 min (System Overview):} Explique el sistema completo desde el NTC hasta la salida del Twin-T[cite: 5].
    \item \textbf{2–5 min (Analytical Probe):} Pregunta Nivel B (requiere dibujar/escribir)[cite: 5].
    \item \textbf{5–8 min (Evidence Probe):} Pregunta Nivel C sobre sus archivos reales[cite: 5].
    \item \textbf{8–11 min (System/Loading):} Pregunta obligatoria sobre loading o $K_L$[cite: 5].
    \item \textbf{11–13 min (Transfer):} Pregunta Nivel D[cite: 5].
\end{enumerate}

\section*{Banco de Preguntas[cite: 5]}
\textbf{Nivel A — Fundamentos[cite: 5]}
\begin{itemize}
    \item \textbf{A3:} ¿Por qué la temperatura modifica principalmente la amplitud de la señal del puente y no la frecuencia de la portadora?[cite: 5]
    \item \textbf{A6:} ¿Por qué la tensión diferencial del puente no debe confundirse con la tensión de uno de sus nodos respecto de tierra?[cite: 5]
\end{itemize}

\textbf{Nivel B — Razonamiento Analítico[cite: 5]}
\begin{itemize}
    \item \textbf{B3:} Explique matemáticamente por qué $V_{dif} = V_{CH1} - V_{CH2}$ y el error de confundirlo con $V_{CH1}$[cite: 5].
    \item \textbf{B5 / B6:} Escriba $V_L = V_{th} \frac{Z_{in}}{Z_{out}+Z_{in}}$ y explique físicamente $K_L \approx 1$ frente a un valor menor[cite: 5].
\end{itemize}

\textbf{Nivel C — Evidencia Real (sobre capturas del estudiante)[cite: 5]}
\begin{itemize}
    \item \textbf{C1 / C3:} Muestre una captura del osciloscopio o Bode. Identifique CH1, CH2, MATH, amplitud, y frecuencia de notch[cite: 5].
    \item \textbf{C5:} Explique cómo transformó un $\Delta t$ en desfase en Python mediante $\phi = 360^\circ f \Delta t$[cite: 5].
\end{itemize}

\textbf{Nivel D — Transferencia[cite: 5]}
\begin{itemize}
    \item \textbf{D2:} Si $Z_{in} \rightarrow \infty$, ¿qué ocurre idealmente con la tensión de la etapa anterior?[cite: 5]
    \item \textbf{D4:} ¿Cómo distinguiría experimentalmente entre loading, un error de conexión, o un error de escala del osciloscopio?[cite: 5]
\end{itemize}

\end{document}
